\section{Công nghệ phần mềm áp dụng trong đồ án}

Giáo trình công nghệ phần mềm~\cite{thanh2024} nhấn mạnh phần mềm là sản phẩm
có vòng đời: xác định nhu cầu, đặc tả, thiết kế, hiện thực, kiểm thử, vận hành.
Đồ án ngành không đủ thời gian cho quy trình thác nước đầy đủ tài liệu IEEE, nhưng
các pha vẫn hiện diện và được ghi trong báo cáo này.

\subsection{Yêu cầu phần mềm}

Yêu cầu chức năng mô tả việc hệ thống phải làm (hỏi đáp, tra cứu, lịch). Yêu cầu
phi chức năng mô tả ràng buộc (thời gian đáp ứng tra cứu, suy giảm khi thiếu
LLM, không CORS mở). Use case là hình thức đặc tả gần người dùng; bảng YC-xx là
hình thức gần kiểm thử. Hai hình thức phải khớp: nếu use case nói 404 khi xem
chéo hội thoại thì phải có test tương ứng.

Yêu cầu thay đổi trong quá trình làm (thêm \texttt{seq}, đổi cổng Postgres) được
ghi như quyết định thiết kế, không giấu. Đó là vòng lặp ngắn, không phải thiếu
phân tích.

\subsection{Thiết kế kiến trúc và thiết kế chi tiết}

Thiết kế kiến trúc chọn phân lớp, chọn client--server, chọn ba container. Thiết
kế chi tiết chọn schema bảng, chọn bảy nút, chọn mã HTTP. Nhầm hai mức này dẫn
tới ``vẽ microservices'' trong khi đồ án một backend --- báo cáo không vẽ kiến
trúc không hiện thực.

Mô hình quan hệ thực thể (ER): thực thể (tổ chức, user, hội thoại, Điều), thuộc
tính, mối kết hợp 1--n. Chuẩn hoá: không lặp nguyên văn luật trong mọi task tuân
thủ; task trỏ rule, rule chứa \texttt{legal\_refs}. Vector không nằm Postgres vì
khác mô hình truy vấn --- đó vẫn là thiết kế, không phải thiếu ER.

\subsection{Hiện thực và chuẩn mã}

Quy ước: Python type hint ở biên giới API; TypeScript strict phía client; không
commit \texttt{.env}. Tên nút LangGraph giữ tiếng Anh khớp mã; báo cáo giải thích
tiếng Việt. Tách router / service / model tránh SQL trong JSX và tránh JSX trong
service.

\subsection{Vận hành và bảo trì}

Compose, healthcheck, volume, Alembic là vận hành mức đồ án. Bảo trì corpus
(import, reindex) nằm trang admin, không đòi sửa mã khi thêm văn bản cùng schema.
Bảo trì rule tuân thủ vẫn đòi sửa Python/seed --- hạn chế đã nêu.

\subsection{Vì sao không microservices}

Tách auth, RAG, FTS thành nhiều dịch vụ tăng độ phức tạp mạng, quan sát và triển
khai, không tăng giá trị với một người hiện thực. Ba container đã tách được mối
quan tâm lưu trữ / API / tĩnh. Hướng sau này: tách worker soát xét và tách Chroma
khi cần nhân bản backend, không phải tách từ ngày đầu.

\subsection{Máy ảo, container và WSL2}

Máy ảo (VM) ảo hoá phần cứng, chạy kernel riêng, nặng. Container chia kernel
máy chủ, đóng gói process và phụ thuộc, nhẹ hơn~\cite{docker}. Đồ án dùng
container. WSL2 là máy ảo Linux trên Windows, là môi trường \emph{host} khi phát
triển, không phải thành phần sản phẩm. Người chấm trên Linux thuần vẫn chạy
cùng file Compose.

\subsection{Linux, tiến trình và cổng}

Mỗi container là một (hoặc vài) tiến trình nghe cổng. Xung đột cổng 5432 trên
máy có sẵn Postgres là lý do đổi 23432. \texttt{localhost} trong container không
phải localhost host --- đã nêu ở mục mạng cầu. Hiểu tiến trình/cổng giúp đọc
\texttt{docker compose ps} khi demo hỏng.
